% =====================================================================
%  05-verificacion-independiente.tex - Verificacion independiente
% =====================================================================

\section{Verificación independiente}\label{sec:pruebas-interop}

Un principio transversal de la suite es que ningún resultado criptográfico se acepta únicamente porque el software propio lo apruebe: cada artefacto que el prototipo produce se verifica también con herramientas independientes, y cada artefacto que las herramientas independientes producen se verifica con el software propio. La herramienta de referencia es OpenSSL, en la versión 3.5.7 (9 de junio de 2026) del entorno de pruebas; la versión queda registrada en tres lugares: en las pruebas (que la leen con \texttt{openssl version}), en el instructivo de cada paquete de evidencia exportado, y en la transcripción de la demostración. Este cruce sistemático materializa la mitigación comprometida en el análisis de riesgos para la discrepancia de digests ante el sellado de tiempo (RLC-01, \cref{tab:rlc-01}), que exigía auditar los resúmenes generados en Rust contra herramientas independientes de línea de comandos. La \cref{tab:pruebas-interop} resume las comprobaciones cruzadas que la suite ejecuta como subprocesos en cada corrida.

\begin{table}[H]
  \centering
  \caption{Comprobaciones de interoperabilidad ejecutadas por la suite con herramientas independientes.}
  \label{tab:pruebas-interop}
  \begin{tabularx}{\linewidth}{@{}L{3.6cm} X@{}}
    \toprule
    \textbf{Artefacto} & \textbf{Comprobación cruzada} \\
    \midrule
    Resumen SHA-256 & El digest del hasher propio coincide con la salida de \texttt{sha256sum} del sistema sobre un archivo de 123\,457 bytes. \\
    \addlinespace
    Firma propia & \texttt{openssl dgst -sha256 -verify} acepta las firmas del prototipo (\texttt{Verified OK}); se ensaya en integración, de extremo a extremo tras \texttt{sign}, y sobre el paquete de evidencia. \\
    \addlinespace
    Firma ajena & Una firma generada con \texttt{openssl dgst -sha256 -sign} es aceptada por el verificador propio: la interoperabilidad es bidireccional. \\
    \addlinespace
    Certificados y CRL & \texttt{openssl verify -crl\_check} concuerda con el validador propio en el certificado vigente y en el revocado; \texttt{openssl verify -attime} concuerda en la cadena vigente y en el emisor expirado. \\
    \addlinespace
    Sellos de tiempo & \texttt{openssl ts -verify} acepta los tokens de la autoridad local, verificados por digest y por documento; \texttt{openssl ts -reply -text} se usa para leer los números de serie emitidos. \\
    \addlinespace
    Paquete ZIP & \texttt{unzip -t} acepta el paquete y la extracción restaura cada byte, incluida una entrada con los 256 valores de byte; con un byte alterado, \texttt{unzip -t} falla nombrando el CRC. \\
    \addlinespace
    Código TOTP & En la demostración, el código se recomputa con \texttt{base32}, \texttt{od} y \texttt{openssl dgst -sha1 -mac HMAC} (sin el software propio) y el verificador del prototipo lo acepta. \\
    \addlinespace
    Instructivo del paquete & La prueba de extremo a extremo del paquete extrae los comandos de \texttt{INSTRUCCIONES.md} y los ejecuta literalmente: las tres comprobaciones pasan sobre el paquete íntegro y rechazan el miembro corrompido. \\
    \bottomrule
  \end{tabularx}
\end{table}

La evidencia operativa de este enfoque comprende dos demostraciones de extremo a extremo. \texttt{scripts/\allowbreak demo.sh} (\cref{sec:impl-cli}), cuya transcripción íntegra, capturada el 12 de julio de 2026 sobre OpenSSL 3.5.7 y re-ejecutada sin errores con posterioridad, está versionada en \texttt{docs/demo-transcript.md}, recorre la CLI. \texttt{scripts/\allowbreak api-demo.sh} levanta PostgreSQL, Redis y la aplicación en puertos efímeros; prueba TOTP, RBAC, revocación, recovery de un solo uso, carga, sellado, verificación, exportación y auditoría; además comprueba que la contraseña y el token no se almacenan en claro y que el vault no contiene el texto original y valida el ZIP con OpenSSL y \texttt{unzip}; las variables de Cincel se eliminan expresamente del entorno. Ambos guiones terminaron con código cero en el corte del 12 de septiembre; las condiciones de autorización, concurrencia y restauración se detallan en la \cref{sec:pruebas-transacciones}. Cada condición inesperada aborta la ejecución. Los extractos siguientes son evidencia histórica de la primera transcripción y no salidas nuevas de ese corte; se abrevian las rutas del directorio temporal y del binario, sin alterar ninguna salida. El primero es el reporte de la verificación integral con los cuatro componentes en verde:

\begin{lstlisting}[language={},numbers=none,caption={Verificación integral con evidencia íntegra (extracto de \texttt{docs/demo-transcript.md}).},label={lst:pruebas-verify}]
+ despacho-cli verify document.txt --sig document.txt.sig
    --tsr document.txt.tsr --cert socia-demo.crt.pem
    --ca ca.crt.pem --crl crl.pem
document:    document.txt
digest:      fb5c0a6097be9246860fa82c89db3576e8a4ba7e...
integrity:   passed: recomputed digest matches the digest
             the timestamp token attests
signature:   passed: signature verifies over the document
             digest with the signer certificate
certificate: passed: status at the evaluation time: valid
             (revocation checked against the supplied list)
timestamp:   passed: token accepted under the trust anchor;
             generated at 2026-07-12T10:19:41Z
verdict:     valid
\end{lstlisting}

El segundo es la verificación del paquete de evidencia exportado usando únicamente OpenSSL, sin el software que lo generó (el mismo procedimiento que seguiría un perito o la contraparte procesal, reproducido paso a paso en el Anexo~\ref{anexo:verificacion}):

\begin{lstlisting}[language={},numbers=none,caption={Verificación del paquete de evidencia con solo OpenSSL (extracto de \texttt{docs/demo-transcript.md}).},label={lst:pruebas-openssl}]
+ openssl dgst -sha256 document.txt
SHA2-256(document.txt)= fb5c0a6097be9246860fa82c89db3576...
+ openssl x509 -in certificado.pem -pubkey -noout
    -out firmante.pub.pem
+ openssl dgst -sha256 -verify firmante.pub.pem
    -signature document.txt.sig document.txt
Verified OK
+ openssl verify -crl_check -CAfile ca-y-crl.pem certificado.pem
certificado.pem: OK
+ openssl ts -verify -data document.txt -in document.txt.tsr
    -CAfile ca.pem
Using configuration from /etc/pki/tls/openssl.cnf
Verification: OK
\end{lstlisting}

El tercero es la manipulación de la bitácora encadenada: la sustitución de una acción por otra de la misma longitud (de modo que solo el encadenamiento criptográfico puede notarlo) y su detección en el índice exacto:

\begin{lstlisting}[language={},numbers=none,caption={Detección de la manipulación de la bitácora (extracto de \texttt{docs/demo-transcript.md}).},label={lst:pruebas-audit}]
+ despacho-cli audit verify-chain
audit chain valid (2 entries)
+ sed -i s/demo.sello/demo.XXXXX/ audit-log.jsonl
+ (expected to fail) despacho-cli audit verify-chain
Error: audit chain broken at index 1
OK: command was rejected, as expected
\end{lstlisting}

La misma transcripción registra la revocación del certificado del firmante y su efecto sobre la verificación: el componente de certificado falla reportando \texttt{revoked (serial 1000)} y el reporte añade que ese estado habla solo del instante de evaluación y no invalida por sí mismo una firma producida durante la vigencia. Esa es la lectura jurídica cuidadosa que el diseño exigió y que las pruebas fijan textualmente en tres niveles: caso de uso, interfaz de línea de comandos y demostración.
